\documentclass[11pt,a4paper]{article}
\usepackage[utf8]{inputenc}
\usepackage[T1]{fontenc}
\usepackage[margin=2.5cm]{geometry}
\usepackage{booktabs}
\usepackage{enumitem}
\usepackage{xcolor}
\usepackage{hyperref}

\definecolor{todo}{RGB}{200,0,0}
\definecolor{done}{RGB}{0,140,0}
\definecolor{info}{RGB}{0,0,180}
\newcommand{\TODO}[1]{\textcolor{todo}{\textbf{TODO:} #1}}
\newcommand{\DONE}[1]{\textcolor{done}{\textbf{DONE:} #1}}
\newcommand{\NOTE}[1]{\textcolor{info}{\textbf{Note:} #1}}

\title{Notes on the Paper:\\
\large Benchmarking ML Approaches for UV Absorption Wavelength Prediction}
\author{Internal Reference}
\date{Last updated: March 4, 2026}

\begin{document}
\maketitle
\tableofcontents
\newpage

% ======================================================================
\section{Paper Narrative: ``Complementary Strengths''}
% ======================================================================

The central argument of the paper is that \textbf{RF and BiGRU have complementary strengths} for UV $\lambda_{\max}$ prediction:

\begin{itemize}
    \item \textbf{RF excels at interpolation}: Best cross-validated performance on the primary dataset (RMSE 31.34 vs.\ 36.45 nm) and all three external datasets. Beats published GCNN and GBFS baselines.
    \item \textbf{BiGRU excels at extrapolation}: Outperforms RF on 16 novel wetlab compounds (MAE 28.6 vs.\ 38.5 nm). Learns transferable sequential patterns beyond fixed substructural features.
    \item \textbf{Practical guidance}: Use RF for screening within known chemical space; use BiGRU (or DL more broadly) when exploring novel molecular architectures.
\end{itemize}

\NOTE{This narrative is stronger than ``RF beats everything'' because it provides nuanced, actionable guidance. Reviewers value honesty and the OOD finding is genuinely novel---most benchmark papers never test true out-of-distribution generalization.}


% ======================================================================
\section{Git History / Fallback Points}
% ======================================================================

\begin{tabular}{lll}
\toprule
\textbf{Hash} & \textbf{Date} & \textbf{Description} \\
\midrule
\texttt{ca3fd27} & 2026-03-04 & Pre-revision: all experiments, tuned RF, original ``RF wins'' framing \\
\texttt{8429a3f} & 2026-03-04 & ``Complementary strengths'' narrative revision \\
\bottomrule
\end{tabular}

\medskip
To revert to pre-revision: \texttt{git checkout ca3fd27 -- benchmark\_paper.tex}


% ======================================================================
\section{Key Results Summary}
% ======================================================================

\subsection{Primary Dataset (Joung+Beard, 18,755 samples)}

Stratified 5-fold CV, 64/16/20 train/val/test split.

\begin{tabular}{lcccc}
\toprule
\textbf{Model} & \textbf{RMSE (nm)} & \textbf{MAE (nm)} & $R^2$ & Pearson $r$ \\
\midrule
RF tuned + Solvent       & \textbf{31.34 $\pm$ 1.82} & \textbf{15.16 $\pm$ 0.44} & 0.914 & 0.956 \\
XGBoost + Solvent        & 33.70 $\pm$ 1.61 & 20.05 $\pm$ 0.48 & 0.900 & 0.950 \\
RF tuned, Solute only    & 34.13 $\pm$ 1.44 & 16.45 $\pm$ 0.33 & 0.900 & 0.950 \\
BiGRU + Solvent          & 36.45 $\pm$ 1.12 & 20.70 $\pm$ 0.51 & 0.884 & 0.940 \\
BiGRU Solute only        & 39.03 $\pm$ 1.37 & 22.53 $\pm$ 0.32 & 0.870 & 0.930 \\
XGBoost (MAE loss)       & 40.75 $\pm$ 1.58 & 21.99 $\pm$ 0.47 & 0.854 & 0.927 \\
\bottomrule
\end{tabular}

\subsection{Cross-Dataset Results}

\begin{tabular}{llccl}
\toprule
\textbf{Dataset} & \textbf{Model} & \textbf{RMSE} & \textbf{Published} & \textbf{Verdict} \\
\midrule
Deep4Chem  & Our RF   & 22.70 & GCNN 26.6  & \textcolor{done}{Beats published!} \\
Deep4Chem  & Our BiGRU & 27.07 & GCNN 26.6  & Close but slightly worse \\
Jung 2024  & Our RF   & 29.82 & GBFS 32.2  & \textcolor{done}{Beats published!} \\
Jung 2024  & Our BiGRU & 36.31 & GBFS 32.2  & Worse than published \\
nablaColors & Our RF  & MAE 34.62 & UniProp MAE 15.97 & 2.2$\times$ gap (3D needed) \\
\bottomrule
\end{tabular}

\subsection{Wetlab Experimental Validation (16 molecules $\times$ 2 solvents)}

\begin{tabular}{lcc}
\toprule
& \textbf{BiGRU} & \textbf{RF (tuned)} \\
\midrule
Overall MAE  & \textbf{28.6 nm} & 38.5 nm \\
Overall RMSE & \textbf{33.9 nm} & 43.7 nm \\
Excl.\ outliers MAE & \textbf{24.2 nm} & 32.6 nm \\
\bottomrule
\end{tabular}

\NOTE{BiGRU beats RF on wetlab data---this is the key finding that drives the ``complementary strengths'' narrative. The two outliers are Cinnamic Acid ($\sim$70 nm error, measurement at edge of scan range) and Octocrylene ($\sim$89 nm RF error, underrepresented cyano-acrylate in training data).}

\subsection{Classification}
Our RF on reconstructed Reaxys: Accuracy 0.879, Sens 0.876, Spec 0.882. Mamede published: Accuracy 0.89. Close match validates methodology.

\subsection{RF Hyperparameter Tuning}
432-config grid search on fold 0 validation set. Best config: $B = 1000$, \texttt{max\_features} = 0.3, radius = 2, $n_{\text{bits}} = 2048$. The critical parameter was \texttt{max\_features} = 0.3 (30\% of features per split) vs.\ default $\sqrt{d}$ ($\sim$1.6\%). All top 10 configs used \texttt{max\_features} = 0.3.


% ======================================================================
\section{Placeholders to Fill In}
% ======================================================================

\begin{enumerate}
    \item \DONE{Table 2 (Solvent Ablation): BiGRU Solute only = RMSE 39.03 $\pm$ 1.37, MAE 22.53 $\pm$ 0.32, $R^2$ = 0.870. Solvent improvement = 7\%.}
    \item \DONE{Abstract updated: ``8\% RMSE reduction for RF, 7\% for BiGRU''.}
\end{enumerate}


% ======================================================================
\section{Revisions Made (What Changed and Why)}
% ======================================================================

\subsection{Complementary Strengths Reframing (March 4)}
\begin{itemize}
    \item \textbf{Abstract}: Changed from ``RF achieves best overall performance'' to ``RF achieves best cross-validated performance'' + added sentence about BiGRU OOD advantage. Ends with practical model selection guidance instead of ``classical ML can match DL''.
    \item \textbf{Introduction contributions}: Added bullet about OOD generalization finding. Changed ``consistently outperforms'' to ``achieves the best cross-validated performance''.
    \item \textbf{Wetlab discussion (Sec 5.7)}: Expanded from 3 paragraphs to 5. Added explanation of \textit{why} ranking reverses (substructure matching vs.\ transferable patterns). Octocrylene as the exemplar of RF failure on novel conjugation. Added concluding paragraph with practical model selection guidance.
    \item \textbf{Conclusion finding \#1}: Changed from ``RF outperforms deep learning'' to ``RF and BiGRU have complementary strengths''. Added summary paragraph after enumerated findings.
    \item \textbf{Conclusion future work}: Added ensemble approach suggestion.
    \item \textbf{Conclusion limitations}: Added note about small wetlab validation set.
\end{itemize}

\subsection{RF Tuning Update (March 4)}
\begin{itemize}
    \item All RF numbers updated from default (RMSE 32.18) to tuned (RMSE 31.34).
    \item RF no-solvent updated from 34.93 to 34.13 (also tuned).
    \item Methods section (Sec 3.3): Now describes grid search methodology and optimal hyperparameters.
    \item 14\% RF vs BiGRU improvement (was 12\%).
    \item Training time updated: 15 min / 100$\times$ (was 12 min / 125$\times$) due to 1000 trees.
    \item Interpretability figures regenerated with tuned model. Feature importance more distributed (top 5 = 26.7\% vs.\ 38.5\% before) because \texttt{max\_features} = 0.3 lets more features contribute.
    \item Wetlab predictions re-run with tuned RF (slightly worse OOD: MAE 38.5 vs.\ 35.8 with default RF).
\end{itemize}

\subsection{Mathematical Model Descriptions (March 4)}
\begin{itemize}
    \item Added Sec 3.3 (RF), 3.4 (XGBoost), expanded 3.5 (BiGRU), added 3.6 (GNN).
    \item Equations: bagging, MSE splitting, Morgan FP, boosting objective, GRU gates, message passing.
    \item Added 5 bib entries: \texttt{breiman2001random}, \texttt{chen2016xgboost}, \texttt{elman1990finding}, \texttt{cho2014learning}, \texttt{gilmer2017neural}.
\end{itemize}

\subsection{Interpretability Section (March 3)}
\begin{itemize}
    \item Added Sec 5.5 with 4 subfigures: feature importance bar chart, solute vs.\ solvent breakdown, cumulative importance curve, training time comparison.
    \item Key stats: solute 96.1\%, solvent 3.9\% (24.6:1 ratio). Top feature (bit 463, aryl diazo) = 13.0\%.
    \item ``Practical advantages'' paragraph uses measured tone per user preference.
\end{itemize}

\subsection{Tone / Language (March 3)}
\begin{itemize}
    \item User requested softer language in ``practical advantages'' paragraph.
    \item Changed from prescriptive (``researchers should evaluate fingerprint-based models as a first approach'') to measured (``depending on the dataset and molecular application domain, it is prudent to evaluate classical ML models where appropriate'').
    \item This applies throughout: never say ``researchers should''; instead frame as findings that inform decisions.
\end{itemize}


% ======================================================================
\section{Potential Reviewer Questions and Prepared Answers}
% ======================================================================

\begin{enumerate}[leftmargin=*]
    \item \textbf{``Is the RF vs BiGRU difference statistically significant?''}\\
    \TODO{Run paired t-tests across 5 folds. With only 5 data points per fold comparison, may not reach significance.} The cross-dataset consistency (RF wins on all 4 datasets) is itself strong evidence.

    \item \textbf{``Why not use more sophisticated DL models (Transformer, GNN)?''}\\
    The benchmark scope is 1D representations (SMILES/fingerprints). GNNs are discussed in context of nablaColors (Sec 5.4) as a 3D comparison point, not as direct competitors. Transformers could be future work.

    \item \textbf{``The wetlab validation set is small (16 molecules).''}\\
    Acknowledged in Limitations. However, the finding that \textit{ranking reverses} on OOD data is qualitatively robust even with small N---it's a directional finding, not a precise quantification.

    \item \textbf{``How were RF hyperparameters selected?''}\\
    Grid search over 432 configurations using fold 0 validation set (Sec 5.1). Best config applied to all 5 folds. XGBoost and BiGRU use standard/published hyperparameters.

    \item \textbf{``Why does the tuned RF perform worse on wetlab molecules?''}\\
    \texttt{max\_features} = 0.3 (vs.\ $\sqrt{d}$) allows the model to use more features per split, which improves fit to the training distribution but may increase overfitting to training-specific substructure patterns. This is consistent with the interpolation/extrapolation distinction.

    \item \textbf{``What about data leakage in v1 results?''}\\
    v1 results (early stopping on test set) are \textit{not used in the paper}. All reported results use v2 protocol (separate validation set for early stopping). v1 is only in supplementary for reference.

    \item \textbf{``Can you provide per-fold breakdowns?''}\\
    \TODO{Add supplementary table with per-fold RMSE, MAE, $R^2$ for all models.}

    \item \textbf{``What about the solvent improvement for BiGRU?''}\\
    \DONE{BiGRU no-solvent v2 complete: RMSE 39.03 $\pm$ 1.37 $\to$ 36.45 $\pm$ 1.12 with solvent = 7\% improvement. Consistent with RF's 8\% improvement, confirming solvent concatenation benefits both model families.}
\end{enumerate}


% ======================================================================
\section{Target Journals}
% ======================================================================

\begin{enumerate}
    \item \textbf{Journal of Cheminformatics} (Springer, IF $\sim$8): Strong fit for benchmark/methodology papers. Good precedent for ``simple baselines beat DL'' papers. Open access.
    \item \textbf{Digital Discovery} (RSC): Newer journal, appreciates practical ML studies.
    \item \textbf{Molecular Informatics} (Wiley, IF $\sim$3): Good backup, more applied focus.
    \item \textbf{J.\ Chem.\ Inf.\ Model.} (ACS, IF $\sim$6): Possible but competitive; may want more novelty.
\end{enumerate}


% ======================================================================
\section{File Reference}
% ======================================================================

\begin{tabular}{ll}
\toprule
\textbf{File} & \textbf{Purpose} \\
\midrule
\texttt{benchmark\_paper.tex} & Main paper (benchmark framing, submit this) \\
\texttt{ml\_chemistry\_template.tex} & Original DL-focused paper (preserved, Overleaf-synced) \\
\texttt{Proposal.bib} & Bibliography (77K tokens, read with offset/limit) \\
\texttt{run\_baselines.py} & Main training script (9 models, 5-fold CV) \\
\texttt{tune\_rf.py} & RF hyperparameter grid search \\
\texttt{eval\_wetlab.py} & Wetlab experimental validation \\
\texttt{eval\_classification.py} & Mamede photosafety classification \\
\texttt{run\_cross\_dataset.py} & Cross-dataset benchmarks \\
\texttt{analyze\_rf\_interpretability.py} & RF feature importance analysis + figures \\
\texttt{results/rf\_best\_config.json} & Tuned RF hyperparameters \\
\texttt{results/wetlab\_predictions.csv} & All wetlab predictions (BiGRU + RF) \\
\texttt{results/rf\_hyperparam\_search.csv} & Full 432-config grid search results \\
\bottomrule
\end{tabular}


% ======================================================================
\section{Checklist Before Submission}
% ======================================================================

\begin{enumerate}
    \item[\DONE{}] Tuned RF results in all tables and text
    \item[\DONE{}] Wetlab predictions updated with tuned RF
    \item[\DONE{}] Interpretability figures regenerated with tuned model
    \item[\DONE{}] ``Complementary strengths'' narrative in abstract, intro, discussion, conclusion
    \item[\DONE{}] Mathematical model descriptions (RF, XGBoost, BiGRU, GNN)
    \item[\DONE{}] Cross-dataset results with published baselines
    \item[\DONE{}] Classification validation
    \item[\DONE{}] Fill in BiGRU no-solvent v2 numbers (Table 2 + abstract)
    \item[\TODO{}] Verify all \texttt{\\cite\{\}} keys exist in \texttt{Proposal.bib}
    \item[\TODO{}] Add supplementary material (per-fold tables)
    \item[\TODO{}] Statistical significance tests (paired t-tests)
    \item[\TODO{}] Proofread entire paper
    \item[\TODO{}] Format for target journal guidelines
    \item[\TODO{}] Add \texttt{.gitignore} for large model files
\end{enumerate}


\end{document}
